% R596: whole EGA I Chapter0 3.5 heading and3.5.1.
% OpenAI Codex, GPT-6.1 Sol, Ultra; no human review claimed.
\subsection{Imágenes inversas de prehaces}

\begin{env}[3.5.1]
\label{0.3.5.1-fr}
Bajo las hipótesis de (3.4.1), si $\mathcal{F}$ (resp. $\mathcal{G}$) es un
prehaz sobre $X$ (resp. $Y$) con valores en $\boldsymbol{K}$, todo morfismo
$u:\mathcal{G}\to\psi_*(\mathcal{F})$ de prehaces sobre $Y$ se llama también
un $\psi$-\emph{morfismo} de $\mathcal{G}$ en $\mathcal{F}$, y se denota
asimismo por $\mathcal{G}\to\mathcal{F}$. Se denota también por
$\operatorname{Hom}_\psi(\mathcal{G},\mathcal{F})$ el conjunto
$\operatorname{Hom}_Y(\mathcal{G},\psi_*(\mathcal{F}))$ de los
$\psi$-morfismos de $\mathcal{G}$ en $\mathcal{F}$.

Para todo par $(U,V)$, donde $U$ es un abierto de $X$ y $V$ un abierto de $Y$
tal que $\psi(U)\subset V$, se tiene un morfismo
$u_{U,V}:\mathcal{G}(V)\to\mathcal{F}(U)$ componiendo el morfismo de
restricción
$\mathcal{F}(\psi^{-1}(V))\to\mathcal{F}(U)$ y el morfismo
$u_V:\mathcal{G}(V)\to\psi_*(\mathcal{F})(V)=\mathcal{F}(\psi^{-1}(V))$~; es
inmediato que estos morfismos hacen conmutativos los diagramas
\begin{equation*}
\begin{tikzcd}[column sep=huge,row sep=large]
\mathcal{G}(V) \arrow[r,"u_{U,V}"] \arrow[d]
  & \mathcal{F}(U) \arrow[d] \\
\mathcal{G}(V') \arrow[r,"u_{U',V'}"']
  & \mathcal{F}(U')
\end{tikzcd}
\tag{3.5.1.1}\label{0.3.5.1.1-fr}
\end{equation*}
para $U'\subset U$, $V'\subset V$, $\psi(U')\subset V'$. Recíprocamente, el
dato de una familia $(u_{U,V})$ de morfismos que hacen conmutativos los
diagramas (3.5.1.1) define un $\psi$-morfismo $u$, pues basta tomar
\[
  u_V=u_{\psi^{-1}(V),V}.
\]

\oldpage[0\textsubscript{I}]{31}

Si la categoría $\boldsymbol{K}$ admite límites proyectivos (generalizados), y
si $\mathfrak{B}$, $\mathfrak{B}'$ son bases de las topologías de $X$ y
$Y$ respectivamente, para definir un $\psi$-morfismo $u$ de \emph{haces},
basta dar los $u_{U,V}$ para
$U\in\mathfrak{B}$, $V\in\mathfrak{B}'$ y $\psi(U)\subset V$, que verifican las
condiciones de compatibilidad (3.5.1.1) para $U$, $U'$ en $\mathfrak{B}$ y
$V$, $V'$ en $\mathfrak{B}'$~; basta, en efecto, definir $u_W$, para
todo abierto $W\subset Y$, como límite proyectivo de los $u_{U,V}$ para
$V\in\mathfrak{B}'$ y $V\subset W$, $U\in\mathfrak{B}$ y
$\psi(U)\subset V$.

Cuando la categoría $\boldsymbol{K}$ admite límites inductivos, se tiene, para
todo $x\in X$, un morfismo
\[
  \mathcal{G}(V)\to\mathcal{F}(\psi^{-1}(V))\to\mathcal{F}_x
\]
para toda vecindad abierta $V$ de $\psi(x)$ en $Y$, y estos morfismos
forman un sistema inductivo que da, por paso al límite, un morfismo
\[
  \mathcal{G}_{\psi(x)}\to\mathcal{F}_x.
\]
\end{env}
